\documentclass[11pt,a4paper]{article}
\usepackage[T1]{fontenc}
\usepackage[utf8]{inputenc}
\usepackage{lmodern}
\usepackage{microtype}
\usepackage[a4paper,margin=1in]{geometry}
\usepackage{xcolor}
\usepackage{booktabs,longtable,tabularx,array}
\usepackage{enumitem}
\usepackage{hyperref}
\usepackage{listings}
\usepackage{fancyhdr}
\usepackage{titlesec}
\usepackage{parskip}
\usepackage{caption}
\usepackage{pifont}

\definecolor{meopblue}{HTML}{143C6B}
\definecolor{meoplight}{HTML}{EAF1F8}
\definecolor{meopmid}{HTML}{6B7C93}
\definecolor{meopgreen}{HTML}{2E7D32}
\definecolor{meopamber}{HTML}{B26A00}
\definecolor{meopred}{HTML}{A3312B}
\definecolor{codebg}{HTML}{F7F8FA}

\hypersetup{
  colorlinks=true,
  linkcolor=meopblue,
  urlcolor=meopblue,
  citecolor=meopblue,
  pdftitle={MEOP Process renovation status and toolbox guide},
  pdfauthor={OpenAI for the MEOP project}
}

\titleformat{\section}{\Large\bfseries\color{meopblue}}{\thesection}{0.7em}{}
\titleformat{\subsection}{\large\bfseries\color{meopblue}}{\thesubsection}{0.7em}{}
\titleformat{\subsubsection}{\normalsize\bfseries\color{meopblue}}{\thesubsubsection}{0.7em}{}

\pagestyle{fancy}
\fancyhf{}
\fancyhead[L]{\textcolor{meopmid}{MEOP Process renovation status}}
\fancyhead[R]{\textcolor{meopmid}{August 22, 2026}}
\fancyfoot[C]{\textcolor{meopmid}{\thepage}}
\setlength{\headheight}{15pt}

\lstdefinestyle{meopcode}{
  basicstyle=\ttfamily\small,
  backgroundcolor=\color{codebg},
  frame=single,
  rulecolor=\color{black!10},
  framerule=0.6pt,
  columns=fullflexible,
  keepspaces=true,
  showstringspaces=false,
  breaklines=true,
  upquote=true
}

\newcommand{\done}{\textcolor{meopgreen}{\ding{51}}}
\newcommand{\partdone}{\textcolor{meopamber}{\ding{108}}}
\newcommand{\todoitem}{\textcolor{meopred}{\ding{55}}}
\newcommand{\pkg}[1]{\texttt{#1}}

\begin{document}

\begin{titlepage}
  \centering
  \vspace*{1.7cm}
  {\Huge\bfseries\color{meopblue}MEOP Process\\[0.25em] Renovation Status and Toolbox Guide\par}
  \vspace{0.8cm}
  {\Large Current software snapshot, operating model, batch workflow, and roadmap\par}
  \vspace{1.2cm}

  \noindent\colorbox{meoplight}{%
    \parbox{0.93\textwidth}{
      \vspace{0.35cm}
      \textbf{Purpose.} This document summarizes what the renovated MEOP toolbox currently does, how it is installed and used, how runtime data are expected to be found, how to run the new resumable multi-deployment workflow, what remains to be validated scientifically across deployments, and the recommended strategy for the next development phases.
      \vspace{0.35cm}
    }
  }

  \vspace{1.2cm}

  \begin{tabularx}{0.93\textwidth}{>{\bfseries}p{0.30\textwidth}X}
    Snapshot date & August 22, 2026 \\
    Current objective & Reach a minimal, functional Python version of the MEOP processing chain that can be re-run over the available deployments with readable logs, resumable state, and updated summary tables \\
    Main regression examples & \pkg{ct88-225-12}, \pkg{ct78-465-12}, and the shortened \pkg{ct96-24-13} full-resolution reference set \\
    Pipeline now covered in Python & Catalog discovery, raw ODV import (including FL sidecar merging), \pkg{lr0}, QC/filtering, \pkg{hr0}, \pkg{hr1/lr1}, \pkg{fr0}, \pkg{fr1}, \pkg{hr2}, standard diagnostics, CORA acquisition/preparation, CORA-based T/S calibration plots and offset diagnostics, and resumable batch processing \\
    Execution model & Single Python execution path; no alternate runtime backend remains in the package \\
    Tests in repository & 223 collected tests in the current snapshot \\
  \end{tabularx}

  \vfill
  {\large Prepared for the MEOP toolbox renovation effort\par}
\end{titlepage}

\tableofcontents
\newpage

\section{Executive summary}

The toolbox has moved from an operator workspace to a Python package with explicit runtime ownership and a nearly complete processing path for the currently targeted products. The package can now bootstrap its runtime layout, discover deployments and tags from CSV and JSON metadata, import raw ODV files, generate the low-resolution and high-resolution profile products, generate full-resolution products from HR text files, produce standard diagnostics, acquire and prepare a traceable CORA reference archive, and run through many deployments in sequence while logging failures without stopping the whole run.

The most important practical change in this phase is operational rather than architectural: there is now a resumable batch runner suitable for reprocessing all available deployments. It keeps a persistent per-deployment state file, skips deployments that already completed successfully only while their canonical outputs still exist unless forced, records one log file per deployment, generates readable Markdown and CSV run reports, reconciles stale successful state entries against the filesystem at startup, and refreshes \pkg{list\_tags.csv} and \pkg{list\_deployments.csv} incrementally at the end of the run.

At this point the project is close to a minimal functional version. The next work is less about scaffolding and more about scientific validation, richer diagnostics, and the remaining publication/reporting helpers.

\section{General description}

\subsection{Scientific intent}

MEOP Process transforms marine mammal CTD observations into Argo-like products and associated diagnostics that can be inspected, adjusted, summarized, and released. The software has to combine heterogeneous metadata sources, low-resolution ODV inputs, optional high-resolution time series, and operator-managed catalog and JSON metadata.

\subsection{Software intent of the renovated toolbox}

The renovated toolbox now has four simultaneous roles:

\begin{enumerate}[leftmargin=1.6em]
  \item a Python package and CLI that can be installed, imported, and scripted in a standard way;
  \item a single Python processing engine for the current workflow;
  \item a runtime data manager that owns where tables, catalogs, raw data, outputs, and batch state live;
\end{enumerate}

\section{Current status of the toolbox}

\subsection{Pipeline ownership}

\begin{longtable}{>{\raggedright\arraybackslash}p{0.31\textwidth} >{\raggedright\arraybackslash}p{0.15\textwidth} >{\raggedright\arraybackslash}p{0.44\textwidth}}
\toprule
\textbf{Workflow element} & \textbf{Owner} & \textbf{Notes} \\
\midrule
\endhead
Configuration loading and path resolution & Python & No external runtime dependency. \\
Runtime data layout and validation & Python & Explicit directories under the canonical \pkg{data/} tree plus the configured \pkg{public/} release tree. \\
Catalog and mirrored JSON metadata resolution & Python & Uses \pkg{list\_deployment.csv}, \pkg{list\_deployment\_hr.csv}, and mirrored deployment/platform JSON files. \\
Raw ODV import and profile indexing & Python & Deployment and tag discovery can be driven directly from raw ODV content; FL sidecar files are automatically merged for CHLA/DOXY. \\
\pkg{create\_ncargo} / \pkg{lr0} & Python & Includes low-resolution netCDF assembly. \\
\pkg{lr0} QC and filtering & Python & Ported to Python. \\
Location adjustment & Python & Applies crawl/CLS/SMRU-based geographic corrections and updates profile metadata. \\
\pkg{hr0} & Python & Regular 1 dbar vertical product. \\
\pkg{apply\_adjustments} & Python & Uses runtime coefficient, parameter, and salinity-offset tables. \\
\pkg{hr1/lr1} no-TLC branch & Python & Implemented and regression-tested. \\
\pkg{hr1/lr1} TLC branch & Python & Uses \pkg{gsw} when available and is regression-tested on reference tags. \\
\pkg{fr0} & Python & Detects ascents from HR text files, associates low-resolution metadata, and writes full-resolution products. \\
\pkg{fr1} & Python & Applies the same adjustment/TLC logic in the FR path. \\
\pkg{hr2} & Python & Uses \pkg{fr1} when FR data exist and falls back to \pkg{hr1} otherwise. \\
Standard diagnostics figures & Python & Per-tag section, T/S, map, profile, and flag panels; per-deployment recap; global overview. \\
CORA source acquisition and tile preparation & Python & \pkg{meop-cora-download} creates a release-pinned CTD/Argo source archive; \pkg{meop-cora-prepare} applies documented type/QC/value policies and writes provenance-preserving tiles and a manifest. \\
CORA-based T/S calibration plots & Python & \pkg{meop-compare --plot1 <smru\_name>} filters invalid/minor target components, checks target support before reference I/O, and loads indexed CORA rows along an anti-meridian-safe track corridor; it writes context plots plus per-tag T/S offset CSV and PNG diagnostics. \pkg{meop-compare-batch} runs method version 4 across deployments. \\
Batch processing and summary-table refresh & Python & Resumable multi-deployment runner with logs and reports. \\
High-end analyst diagnostics and publication redesign & Future work & Still to be designed in depth. \\
\bottomrule
\end{longtable}

\subsection{Installed CORA reference archive}

The configured reference archive is derived from Copernicus Marine product
\path{INSITU_GLO_PHY_TS_DISCRETE_MY_013_001}, dataset version \pkg{202511}. The source
download contains 34,556 NetCDF files selected from raw CORA \pkg{CT/OC/TE} candidates
and EasyCORA \pkg{PF} files. The installed preparation contains 540 populated 10-degree
tiles and 4,328,083 profiles: 1,722,428 CTD profiles selected with
\pkg{PROBE\_TYPE=2} and 2,605,655 Argo profiles selected with \pkg{PROBE\_TYPE=5}.

CTD observations cover 1953-04-09 through 2025-06-30, and Argo observations cover
1995-09-07 through 2025-06-30. All tiles share a 10--1000 dbar pressure grid at 10-dbar
spacing and a common \pkg{JULD} origin. Preparation accepts QC flags 1 and 2, uses an
adjusted whole-profile parameter when available and raw values otherwise, and requires
both temperature and salinity. The manifest reports 400,035 profiles excluded by the
strict time/position policy and 543,873 with insufficient paired T/S support; no source
file was skipped and no expected QC variable was missing.

The source and preparation manifests are stored with the data. The configured tile root
is \path{../REF_DATASETS/CORA_data/CORA_ctd_argo_202511_tiles/}. Structural validation
confirmed the profile totals and type split, tile bounds, coordinates, pressure grid and
date coverage. Calibration outputs must be rebuilt whenever this reference archive or
its preparation policy changes.

\subsection{Reference cases currently used}

The current regression set is built around three representative cases plus four additional fixture cases added to broaden sensor and processing-path coverage:

\begin{itemize}[leftmargin=1.5em]
  \item \pkg{ct88-225-12}: low-resolution and TLC validation without full-resolution data;
  \item \pkg{ct78-465-12}: a second low-resolution reference to increase confidence in the non-FR path;
  \item \pkg{ct96-24-13}: shortened full-resolution reference set used to validate \pkg{fr0}, \pkg{fr1}, \pkg{traj}, and \pkg{hr2};
  \item \pkg{ct160-Oxy726-20}: DOXY via FL sidecar ODV (\pkg{OXY\_CONT\_20A} firmware);
  \item \pkg{ft11-Cy07b-12}: CHLA via FL sidecar ODV (\pkg{FTD\_10A} firmware), no DOXY;
  \item \pkg{ct153-Pendragon-19}: plain CTD-only tag, no full-resolution, tests LR-only \pkg{apply\_adjustments};
  \item \pkg{ct107-938-13-N2}: split-tag fixture (base \pkg{ct107-938-13} split into N1/N2 by \pkg{table\_split\_tags.csv}).
\end{itemize}

These cases are valuable because they test both the low-resolution/high-resolution workflow and the branch that uses full-resolution HR text inputs.

\section{What has been completed}

\subsection{Repository and package structure}

The project now uses a standard \pkg{src/} layout with \pkg{pyproject.toml}, a package root at \pkg{src/meop\_process/}, a public API module, a CLI, repository wrapper scripts under \pkg{scripts/}, and deprecated historical code under \pkg{legacy/}.

\subsection{Runtime data ownership}

The package now owns its runtime layout explicitly. In practice this means that the code no longer relies on implicit knowledge of where files are stored. It has named roots for:

\begin{itemize}[leftmargin=1.5em]
  \item packaged CSV defaults;
  \item operator-maintained catalogs;
  \item mirrored JSON metadata;
  \item raw ODV inputs;
  \item raw HR CTD inputs;
  \item processed profile and trajectory outputs;
  \item plots, diagnostics, and maps;
  \item batch logs, reports, and resumable state.
\end{itemize}

\subsection{Scientific workflow stages already ported}

The following steps are already owned by Python and exercised in tests:

\begin{itemize}[leftmargin=1.5em]
  \item deployment and tag resolution;
  \item raw ODV import and parsing;
  \item \pkg{lr0};
  \item the low-resolution QC/filter block;
  \item \pkg{hr0};
  \item \pkg{apply\_adjustments};
  \item \pkg{hr1} and \pkg{lr1} for both no-TLC and TLC branches;
  \item \pkg{fr0};
  \item \pkg{fr1};
  \item \pkg{hr2};
  \item standard diagnostics figures;
  \item CORA source inventory/download and CTD/Argo tile preparation;
  \item CORA-based T/S calibration plots and offset diagnostics
  (\pkg{meop-compare --plot1} and \pkg{meop-compare-batch}).
\end{itemize}

\subsection{Operational batch runner}

A dedicated multi-deployment runner is now present. Its design goals are:

\begin{itemize}[leftmargin=1.5em]
  \item continue past deployment-level failures;
  \item make reruns cheap by skipping deployments that already completed successfully;
  \item produce one log file per deployment;
  \item produce a run-level Markdown summary and CSV table;
  \item reconcile stale successful state entries when outputs have been deleted;
  \item refresh summary metadata tables at the end without reopening every netCDF file when nothing changed.
\end{itemize}

\section{Installation}

\subsection{Editable install}

The recommended installation uses the repository's \pkg{uv} environment:

\begin{lstlisting}[style=meopcode]
uv sync --extra dev --extra cora
\end{lstlisting}

\subsection{Expected Python dependencies}

The current minimal scientific runtime expects at least:

\begin{itemize}[leftmargin=1.5em]
  \item \pkg{numpy}
  \item \pkg{pandas}
  \item \pkg{xarray}
  \item \pkg{scipy}
  \item \pkg{h5netcdf}
  \item \pkg{h5py}
  \item \pkg{matplotlib}
  \item \pkg{gsw}
\end{itemize}

If map backgrounds are desired in diagnostics, \pkg{cartopy} can be installed as an optional extra. The code currently falls back to a plain longitude/latitude plot when \pkg{cartopy} is unavailable.

\section{Runtime data layout}

\subsection{Canonical runtime roots}

\begin{itemize}[leftmargin=1.5em]
  \item \textbf{Packaged/default tables:} \pkg{src/meop\_process/resources/tables/}, mirrored into \pkg{data/tables/}.
  \item \textbf{Operator-maintained catalogs:} \pkg{data/catalog/}, notably \pkg{list\_deployment.csv} and \pkg{list\_deployment\_hr.csv}.
  \item \textbf{Mirrored metadata JSON:} \pkg{data/data\_raw/config\_files/}, for example \pkg{deployment2.json}, \pkg{platform2.json}, and patch files.
  \item \textbf{Low-resolution raw inputs:} \pkg{data/data\_raw/raw\_smru\_data\_odv/}.
  \item \textbf{High-resolution raw inputs:} \pkg{data/data\_raw/raw\_smru\_hr\_data/<year>/}, containing files such as \pkg{<prefix><instr\_id>\_ctd.txt}.
  \item \textbf{Processed profile outputs:} \pkg{data/data\_prof/}, containing \pkg{lr0}, \pkg{lr1}, \pkg{hr0}, \pkg{hr1}, \pkg{hr2}, \pkg{fr0}, and \pkg{fr1} profile products.
  \item \textbf{Processed trajectory outputs:} \pkg{data/data\_traj/}.
  \item \textbf{CORA 202511 source archive:} \path{../REF_DATASETS/CORA_data/CORA_ctd_argo_source_202511/}.
  \item \textbf{Prepared CORA tiles:} \path{../REF_DATASETS/CORA_data/CORA_ctd_argo_202511_tiles/}.
  \item \textbf{Diagnostics by tag:} \pkg{data/plots\_by\_tags/}.
  \item \textbf{Deployment and overview plots:} \pkg{data/plots\_by\_deployments/} and \pkg{data/plots\_overview/}.
  \item \textbf{Maps:} \pkg{data/maps/}.
  \item \textbf{Batch state and reports:} \pkg{data/batch/}.
\end{itemize}

\subsection{How raw data are expected to be found}

Low-resolution data are staged directly under \pkg{data/data\_raw/raw\_smru\_data\_odv/}. High-resolution files are resolved from \pkg{list\_deployment\_hr.csv} and are expected under the runtime HR tree, typically as files named like \pkg{<optional\_prefix><instr\_id>\_ctd.txt} inside the appropriate year directory.

This is now explicit in Python and does not rely on any alternate runtime.

\section{Usage}

\subsection{Single deployment processing}

Installed entry point:

\begin{lstlisting}[style=meopcode]
meop-process --deployment ct88 --process_data --diagnostics
\end{lstlisting}

Repository wrapper:

\begin{lstlisting}[style=meopcode]
python scripts/process_deployment.py --deployment ct88 --process_data --diagnostics
\end{lstlisting}

\subsection{Batch processing over all deployments}

Installed entry point:

\begin{lstlisting}[style=meopcode]
meop-process --run-all-deployments
\end{lstlisting}

Dedicated batch entry point:

\begin{lstlisting}[style=meopcode]
meop-process-batch
\end{lstlisting}

Repository wrapper:

\begin{lstlisting}[style=meopcode]
python scripts/run_all_deployments.py
\end{lstlisting}

Useful variants:

\begin{lstlisting}[style=meopcode]
meop-process --run-all-deployments
meop-process-batch
meop-process --run-all-deployments --diagnostics
meop-process --run-all-deployments --no-diagnostics
meop-process --run-all-deployments --force-failed
meop-process --run-all-deployments --force
meop-process --run-all-deployments --deployment ct96
meop-process --run-all-deployments --jobs 8 --verbose
python scripts/run_all_deployments.py --diagnostics
python scripts/run_all_deployments.py --force-failed
python scripts/run_all_deployments.py --force
python scripts/run_all_deployments.py --deployment ct96
python scripts/run_all_deployments.py --jobs 8 --verbose
\end{lstlisting}

By default, batch diagnostics follow \pkg{defaults.batch.diagnostics} from root \pkg{configs.json} (set to true in the generated bootstrap template). The command-line \pkg{--no-diagnostics} flag disables diagnostics explicitly for one run.

\subsection{What the batch runner writes}

By default the runner writes under \pkg{data/batch/}:

\begin{itemize}[leftmargin=1.5em]
  \item \pkg{latest/deployment\_status.json}: persistent deployment-level state;
  \item \pkg{runs/<timestamp>/logs/<deployment>.log}: one log per deployment;
  \item \pkg{runs/<timestamp>/summary.md}: human-readable run report;
  \item \pkg{runs/<timestamp>/summary.csv}: machine-readable run table.
\end{itemize}

A deployment marked successful is skipped on the next batch run only while its canonical outputs still exist, unless a force option is used. At batch startup, stale successful entries whose outputs have disappeared are pruned from \pkg{latest/deployment\_status.json} before skip decisions are made.

\subsection{Clean CORA calibration refresh}

The configured \pkg{references.cora\_dir} points to the CORA 202511 tiles. A reference
change requires regeneration of both the calibration figures and the per-tag offset
diagnostics. The cleanup below is deliberately restricted to calibration PNG/CSV files
and the calibration batch state; it leaves standard diagnostics and processed profile
products untouched.

\begin{lstlisting}[style=meopcode]
cd /media/disk2/roquet/MEOP_process

uv run meop-compare-batch --dry-run --include-disabled

find data/plots_by_tags -type f \
  \( -name '*_calibration*.png' -o \
     -name '*_calibration_offsets.csv' \) \
  -print -delete

find data/batch/compare -depth -mindepth 1 -print -delete

uv run meop-compare-batch --include-disabled --force -j 2
\end{lstlisting}

The unrestricted catalog selection currently contains 2,717 processed tags; without
\pkg{--include-disabled}, 2,665 tags are selected. The batch report separates successful
results from \pkg{no\_reference}, \pkg{insufficient\_target},
\pkg{insufficient\_reference}, \pkg{invalid\_target} and unexpected \pkg{failed}
outcomes. Method version 4 filters invalid/null-island and minor disconnected target
components, checks support before CORA access, and loads indexed reference rows only
from the union of buffered track cells. It classifies a result as successful only when both temperature and
salinity have at least 10 contributing target profiles, at least three valid levels in
400--600 dbar, finite band and linear estimates, and observed coverage spanning at least
200--600 dbar. Incompatible success state is recomputed. Each successful tag has
one \pkg{\_calibration\_offsets.csv}, one \pkg{\_calibration\_offsets.png}, and one or
more chunked calibration context plots under \pkg{data/plots\_by\_tags/<deployment>/}.

A read-only preflight benchmark over all 2,717 target files on August 22, 2026 classified
1,220 as capable of meeting the target-only support gate, 1,386 as
\pkg{insufficient\_target}, and 111 as \pkg{invalid\_target} because no valid deployment
position remained. Eligible tracks requested a median of 10.5 tiles (90th percentile 19;
maximum 31), instead of every tile in a track-spanning rectangle. On the current host, a
dense-tile test selecting three reference profiles took 1.7 seconds and approximately
501 MiB peak resident memory with indexed row loading, compared with 4.0 seconds and
approximately 684 MiB when the full tile arrays were materialised before masking. These
figures describe this archive and host; they are a regression baseline, not a general
runtime guarantee.

\subsection{Runtime configuration policy}

Runtime configuration is loaded only from:

\begin{itemize}[leftmargin=1.5em]
  \item explicit \pkg{--config-file};
  \item \pkg{MEOP\_CONFIG\_FILE} environment variable;
  \item root \pkg{configs.json} under the process directory.
\end{itemize}

Legacy fallback to \pkg{data/configs.json} is intentionally disabled. Machine selection order is: \pkg{--machine}, then \pkg{MEOP\_MACHINE}, then \pkg{defaults.machine}, then automatic key detection.

\subsection{Metadata summary tables}

At the end of a batch run the package refreshes \pkg{list\_tags.csv} and \pkg{list\_deployments.csv}. The refresh is incremental:

\begin{itemize}[leftmargin=1.5em]
  \item deployments processed in the current run are refreshed;
  \item deployments whose tag inventory changed are refreshed;
  \item unchanged deployments are left untouched.
\end{itemize}

This makes reruns much cheaper once the main outputs already exist.

\section{Test cases and validation strategy}

\subsection{Current regression philosophy}

The current test suite is built around a small number of carefully chosen reference cases rather than a full global comparison across every historical product. The goal has been to make the Python workflow operational first, then broaden scientific validation deployment by deployment.

\subsection{Current test coverage}

The repository currently contains 223 collected tests covering:

\begin{itemize}[leftmargin=1.5em]
  \item configuration loading and data-layout bootstrap;
  \item raw ODV import and catalog resolution;
  \item \pkg{ct88} \pkg{lr0/lr1/hr0/hr1} regressions;
  \item \pkg{ct78} low-resolution and no-FR \pkg{hr2} regression;
  \item shortened \pkg{ct96} \pkg{fr0/fr1/hr2/traj} regression;
  \item \pkg{ct160} DOXY via FL sidecar ODV;
  \item \pkg{ft11} CHLA via FL sidecar ODV;
  \item \pkg{ct153} LR-only \pkg{apply\_adjustments} path (no \pkg{hr2});
  \item \pkg{ct107} split-tag processing;
  \item CORA tile discovery and loading, including padded/unpadded tile names and
  bounding-box filtering;
  \item CORA source inventory selection, manifest handling, CTD/Argo preparation,
  profile typing, QC, interpolation and provenance;
  \item CORA-based T/S calibration plot generation, invalid/minor target-position
  filtering, occupied-cell corridors, indexed reference-row loading, early target
  rejection, offset CSV/PNG diagnostics, chunking, single-tag CLI and resumable
  calibration batch statuses;
  \item diagnostics output generation;
  \item metadata summary refresh;
  \item resumable batch-run behavior.
\end{itemize}

\subsection{What still needs validation}

The next validation phase should use the real local environment to compare the Python outputs against existing reference products over a wider set of deployments. The important point is that the current package is now structured so that such comparisons can be made systematically rather than manually.

\section{API overview}

The stable public surface is intentionally small and centered in \pkg{src/meop\_process/api.py}. The most relevant functions at the current stage are:

\begin{itemize}[leftmargin=1.5em]
  \item \pkg{bootstrap\_data\_store}
  \item \pkg{describe\_runtime\_data\_layout}
  \item \pkg{validate\_runtime\_data\_layout}
  \item \pkg{load\_info\_deployment}
  \item \pkg{import\_raw\_data}
  \item \pkg{process\_tags}
  \item \pkg{apply\_adjustments}
  \item \pkg{create\_fr0}
  \item \pkg{create\_hr2}
  \item \pkg{generate\_diagnostics}
  \item \pkg{update\_metadata\_summaries}
  \item \pkg{run\_all\_deployments}
\end{itemize}

The design intent is to keep this surface stable while allowing the implementation below it to continue evolving.

\section{What remains to be done}

\subsection{Scientific and operational work still open}

The main remaining items are now concentrated in a few areas:

\begin{itemize}[leftmargin=1.5em]
  \item broader scientific validation over many deployments;
  \item richer analyst-facing comparison tooling and integration of additional reference climatologies;
  \item redesign of the publication/reporting helpers;
  \item broader scientific robustness work for remaining problematic deployments.
\end{itemize}

\subsection{Why diagnostics remain important}

Although a standard Python diagnostics set now exists and a CORA-based T/S calibration plot layer is operational, the future diagnostic layer is still a major work package. It should eventually help an analyst inspect profiles, compare them with reference datasets, understand the consequences of proposed adjustments, and archive the reasoning used during delayed-mode processing.

\section{Planned strategy for future developments}

The recommended next steps are:

\begin{enumerate}[leftmargin=1.6em]
  \item use the new batch runner to process as many deployments as practical in the real environment;
  \item review the per-deployment logs and batch summary to identify systematic failures;
  \item compare Python outputs with reference outputs on a representative set of deployments;
  \item strengthen the metadata summaries and diagnostics where operational review shows they are still too light;
  \item redesign the reporting and publication layer in Python once the scientific core is trusted.
\end{enumerate}

The important change is that the project is now in a position where large-scale testing is practical. Earlier phases were about creating the package and porting key workflow stages. The current phase is about exercising that workflow widely, collecting evidence, and using that evidence to prioritize the next scientific and operational refinements.

\section{Conclusion}

The toolbox is now close to a minimal functional version. It can be installed as a Python package, run over deployments through a single execution path, resume large processing campaigns, produce readable logs and reports, and refresh the main summary tables efficiently. The next work should focus on scientific confidence and analyst workflow quality rather than on basic software scaffolding.

\end{document}
